\documentclass[../DoAn.tex]{subfiles}
\begin{document}

Chương 1 đã trình bày bối cảnh, mục tiêu và định hướng tổng quát của đề tài. Trên cơ sở đó, chương này tập trung làm rõ phạm vi nghiên cứu, ngữ cảnh bài toán, các hướng giải pháp liên quan và những nền tảng kỹ thuật cần thiết để hiểu phương pháp được đề xuất ở chương sau. Trọng tâm của chương không phải là mô tả các chức năng phần mềm, mà là xây dựng lập luận nghiên cứu: vì sao bài toán hỏi đáp trên SGK Tin học THPT cần một hệ thống RAG có khả năng truy xuất theo phạm vi, xếp hạng lại ngữ cảnh và đánh giá riêng từng thành phần trong pipeline \cite{lewis2020rag}.

Nội dung chương được tổ chức thành bốn phần. Mục đầu tiên xác định phạm vi nghiên cứu và các giới hạn của đồ án. Mục thứ hai trình bày tổng quan bài toán hỏi đáp giáo dục dựa trên tài liệu. Mục thứ ba khảo sát các nhóm giải pháp liên quan và chỉ ra khoảng trống còn tồn tại. Mục cuối cùng trình bày các kiến thức nền tảng và hướng tiếp cận mà đồ án lựa chọn.

\section{Phạm vi nghiên cứu}

Đối tượng nghiên cứu của đồ án là hệ thống hỏi đáp dựa trên tài liệu (document-grounded QA) cho môn Tin học THPT. Hệ thống tiếp nhận câu hỏi tự nhiên của người dùng, truy xuất ngữ cảnh liên quan từ SGK, sau đó sử dụng LLM để sinh câu trả lời dựa trên phần ngữ cảnh đã được truy xuất. Cách tổ chức này kế thừa hướng tiếp cận hỏi đáp mở có truy xuất tài liệu và sinh văn bản tăng cường truy xuất \cite{karpukhin2020dpr,lewis2020rag}. Vì vậy, bài toán không được xem là một chatbot hội thoại mở thông thường, mà là một hệ thống hỏi đáp có ràng buộc nguồn tri thức rõ ràng.

Phạm vi dữ liệu của đề tài gồm SGK Tin học lớp 10, 11 và 12 thuộc hai bộ sách Cánh Diều (CD) và Kết Nối Tri Thức (KNTT). Đây là nguồn dữ liệu có cấu trúc phân cấp theo bộ sách, lớp, chủ đề, bài học và mục nội dung. Cấu trúc này có ý nghĩa quan trọng đối với bài toán truy xuất (retrieval), vì nhiều khái niệm Tin học có thể xuất hiện ở nhiều lớp hoặc nhiều bộ sách với mức độ trình bày khác nhau. Nếu hệ thống chỉ tìm kiếm theo toàn bộ kho dữ liệu mà không xét phạm vi, câu trả lời có thể lấy nhầm nội dung thuộc lớp hoặc bộ sách khác.

Đối tượng sử dụng chính của hệ thống gồm học sinh, giáo viên và người tự học môn Tin học THPT. Với học sinh, hệ thống cần hỗ trợ đặt câu hỏi tự nhiên, giải thích khái niệm, ôn tập nội dung bài học và định hướng luyện tập. Với giáo viên, hệ thống có thể hỗ trợ chuẩn bị học liệu như dàn ý slide, giáo án hoặc câu hỏi kiểm tra. Tuy nhiên, các tác vụ sinh học liệu này đều phụ thuộc vào một nền tảng chung là khả năng truy xuất đúng ngữ cảnh từ SGK trước khi sinh nội dung.

Trọng tâm nghiên cứu của đồ án nằm ở luồng xử lý (pipeline) truy xuất và sinh văn bản trong kiến trúc Sinh văn bản tăng cường truy xuất (Retrieval-Augmented Generation - RAG). Thành phần truy xuất được đánh giá định lượng thông qua bộ đánh giá chuẩn (benchmark) truy hồi, trong đó một câu hỏi kiểm thử gắn với một nguồn đúng kỳ vọng. Thành phần sinh văn bản được đánh giá theo hướng chất lượng đầu ra và mức độ bám nguồn, ưu tiên tiêu chí đánh giá (rubric) thủ công hoặc kết quả đánh giá đã có sẵn thay vì mở rộng chi phí đánh giá tự động bằng LLM. Các thành phần như giao diện, API, quản lý phiên, triển khai hạ tầng và đóng gói phần mềm được trình bày ở mức cần thiết để chứng minh hệ thống có thể vận hành và kiểm thử, nhưng không phải trọng tâm nghiên cứu chính.

Trong phạm vi đồ án, các tác vụ phụ như sinh slide, soạn giáo án hoặc tạo câu hỏi luyện tập được xem là các luồng ứng dụng trên cùng nền tảng RAG và đa tác tử (multi-agent). Chúng có vai trò chứng minh khả năng mở rộng của kiến trúc, nhưng phần đánh giá định lượng sâu vẫn tập trung vào truy xuất ngữ cảnh và sinh câu trả lời bám tài liệu. Các hướng mở rộng như truy xuất đồ thị (graph retrieval) hoặc tối ưu DevOps không được xem là đóng góp chính nếu chưa có thực nghiệm đầy đủ.

Từ phạm vi trên, câu hỏi nghiên cứu chính của đồ án có thể phát biểu như sau: làm thế nào để truy xuất đúng ngữ cảnh từ SGK Tin học THPT và sinh câu trả lời bám tài liệu, đồng thời đánh giá được chất lượng của từng bước trong luồng RAG. Câu hỏi này dẫn trực tiếp đến các quyết định thiết kế ở Chương 3, đặc biệt là phân đoạn (chunking) theo cấu trúc SGK, truy hồi lai giữa tìm kiếm từ khóa và tìm kiếm vector, xếp hạng lại (reranking), truy xuất thích ứng (adaptive retrieval) và đánh giá truy xuất theo một nguồn đúng duy nhất (single-gold).

\section{Tổng quan về bài toán}

\subsection{Bài toán chatbot giáo dục}

Việc ứng dụng AI trong giáo dục tạo ra nhu cầu về các hệ thống có khả năng tương tác tự nhiên với người học. Thay vì chỉ tra cứu tài liệu theo từ khóa, người học có xu hướng đặt câu hỏi bằng ngôn ngữ tự nhiên, yêu cầu giải thích lại khái niệm, đưa ví dụ hoặc hỗ trợ ôn tập theo một bài học cụ thể. Với môn Tin học THPT, nhu cầu này càng rõ ràng vì nội dung học bao gồm cả lý thuyết, thao tác thực hành, thuật ngữ chuyên ngành, quy trình xử lý dữ liệu và các khái niệm liên quan đến lập trình.

Một chatbot giáo dục dùng cho SGK cần đáp ứng đồng thời yêu cầu về ngôn ngữ tự nhiên và yêu cầu về tính chuẩn xác theo tài liệu. Nếu chỉ ưu tiên khả năng hội thoại, hệ thống có thể tạo câu trả lời mạch lạc nhưng không đúng phạm vi bài học. Ngược lại, nếu chỉ dựa vào tìm kiếm từ khóa, hệ thống khó xử lý các câu hỏi được diễn đạt lại hoặc có ngữ cảnh hội thoại. Do đó, bài toán đặt ra không chỉ là tạo một giao diện hỏi đáp, mà là thiết kế một pipeline có khả năng truy xuất nguồn phù hợp và sinh câu trả lời có căn cứ.

Dữ liệu SGK Tin học THPT có một số đặc điểm khác với tài liệu văn bản thông thường. Nội dung được tổ chức theo lớp, bộ sách, chủ đề, bài học và mục nhỏ. Một số chủ đề như dữ liệu, thuật toán, mạng máy tính hoặc đạo đức số có thể xuất hiện ở nhiều lớp nhưng với độ sâu khác nhau. Ngoài ra, mỗi bài học thường kết hợp nhiều loại nội dung như khái niệm, ví dụ, hoạt động, bài tập, ghi nhớ hoặc câu hỏi vận dụng. Những đặc điểm này khiến hệ thống cần nhận biết vai trò của từng đoạn tài liệu thay vì xem toàn bộ văn bản như một tập đoạn rời rạc không có cấu trúc.

\subsection{Hệ thống hỏi đáp dựa trên tài liệu}

Hệ thống hỏi đáp có thể được chia thành hai nhóm lớn: hỏi đáp đóng (closed-book QA) và hỏi đáp mở (open-book QA). Trong hỏi đáp đóng, mô hình trả lời chủ yếu dựa trên tri thức đã được học trong quá trình huấn luyện. Cách tiếp cận này phù hợp với tri thức phổ quát, nhưng không bảo đảm chính xác khi câu hỏi phụ thuộc vào một tài liệu cụ thể. Trong hỏi đáp mở hoặc hỏi đáp dựa trên tài liệu, hệ thống được cung cấp thêm tài liệu nguồn và câu trả lời phải dựa trên tài liệu đó.

Đồ án này thuộc nhóm hỏi đáp dựa trên tài liệu. Khi người dùng hỏi một khái niệm trong SGK Tin học, hệ thống không nên chỉ dựa vào tri thức tham số của LLM, vì tri thức này có thể không trùng với chương trình giáo dục Việt Nam hoặc có thể vượt khỏi phạm vi lớp học. Thay vào đó, hệ thống cần tìm các phân đoạn (chunk) liên quan trong SGK, đưa các phân đoạn này vào ngữ cảnh, rồi yêu cầu LLM sinh câu trả lời dựa trên ngữ cảnh được cung cấp.

Luồng xử lý tổng quát của một hệ thống hỏi đáp dựa trên tài liệu gồm tiếp nhận câu hỏi, phân tích ý định hoặc phạm vi, truy xuất tài liệu liên quan, xây dựng ngữ cảnh, sinh câu trả lời và trả kết quả cho người dùng. Trong luồng này, truy xuất là bước có ảnh hưởng trực tiếp đến chất lượng đầu ra. Nếu ngữ cảnh bị thiếu hoặc sai nguồn, LLM có thể sinh câu trả lời không đầy đủ hoặc không đúng dù câu lệnh được thiết kế tốt. Vì vậy, đánh giá riêng khâu truy xuất là yêu cầu quan trọng khi nghiên cứu một hệ thống RAG.

\subsection{Thách thức đặc thù của dữ liệu SGK Tin học}

Thách thức đầu tiên đến từ cấu trúc phân cấp của SGK. Cùng một thuật ngữ có thể xuất hiện trong nhiều bài hoặc nhiều lớp, nhưng ý nghĩa sư phạm và mức độ chi tiết không giống nhau. Ví dụ, một khái niệm liên quan đến dữ liệu có thể được trình bày ở lớp 10 như kiến thức nền, nhưng ở lớp 11 hoặc 12 lại gắn với bài toán xử lý dữ liệu hoặc cơ sở dữ liệu. Vì vậy, hệ thống cần lưu và sử dụng siêu dữ liệu (metadata) về bộ sách, lớp, chủ đề và bài học trong quá trình truy xuất.

Thách thức thứ hai là cách học sinh đặt câu hỏi thường không trùng với câu chữ trong SGK. Người học có thể hỏi bằng ngôn ngữ đời thường, dùng từ đồng nghĩa, hỏi tiếp theo một lượt hội thoại trước đó hoặc bỏ qua thông tin về bài học đang xét. Điều này làm cho tìm kiếm từ khóa đơn thuần chưa đủ. Hệ thống cần kết hợp tìm kiếm vector, viết lại truy vấn (query rewriting) hoặc cơ chế phân tích ngữ cảnh để tăng khả năng tìm đúng tài liệu.

Thách thức thứ ba là tính đa dạng của nội dung Tin học. Các câu hỏi về định nghĩa, quy trình, so sánh, ví dụ và vận dụng cần các loại ngữ cảnh khác nhau. Câu hỏi định nghĩa thường cần đoạn mô tả khái niệm trực tiếp, trong khi câu hỏi quy trình có thể cần nhiều bước nằm trong cùng một mục. Nếu việc phân đoạn làm mất mạch nội dung hoặc truy xuất trả về các đoạn rời rạc, câu trả lời cuối cùng có thể thiếu tính sư phạm.

Những thách thức trên cho thấy yêu cầu cốt lõi của bài toán không chỉ là tích hợp LLM vào hệ thống học tập. Hệ thống cần một nền tảng tri thức được chuẩn hóa, cơ chế truy xuất có xét phạm vi, bước xếp hạng lại để cải thiện thứ tự ứng viên (candidate), và quy trình đánh giá có thể xác định lỗi đến từ truy xuất hay khâu sinh văn bản.

\section{Khảo sát các giải pháp liên quan}

\subsection{Các nhóm giải pháp hiện có}

Các giải pháp liên quan đến bài toán trợ lý học tập có thể được chia thành một số nhóm chính. Nhóm thứ nhất là các chatbot dựa trên luật (rule-based) hoặc hệ thống phân loại ý định (intent classification) truyền thống. Các hệ thống này sử dụng luật, kịch bản hoặc bộ ý định được định nghĩa trước để tạo phản hồi. Ưu điểm của nhóm này là dễ kiểm soát trong phạm vi nhỏ và ít xảy ra hiện tượng ảo giác (hallucination) nếu câu hỏi nằm đúng kịch bản. Tuy nhiên, chúng kém linh hoạt khi người dùng diễn đạt khác với mẫu dự kiến, đồng thời chi phí mở rộng luật tăng nhanh khi số bài học và số dạng câu hỏi tăng lên.

Nhóm thứ hai là các hệ thống FAQ hoặc tìm kiếm từ khóa. Cách tiếp cận này phù hợp với các kho câu hỏi cố định hoặc tài liệu có từ khóa rõ ràng. Trong bối cảnh SGK, tìm kiếm từ khóa có thể xử lý tốt các câu hỏi chứa thuật ngữ chính xác, tên khái niệm hoặc cụm từ gần giống tài liệu. Hạn chế của nhóm này là khả năng hiểu câu hỏi tự nhiên còn thấp, đặc biệt khi học sinh diễn đạt lại bằng ngôn ngữ đời thường hoặc hỏi theo ngữ cảnh hội thoại.

Nhóm thứ ba là tìm kiếm ngữ nghĩa (semantic search) dựa trên vector nhúng (embedding). Thay vì chỉ so khớp từ khóa, tìm kiếm ngữ nghĩa biểu diễn câu hỏi và đoạn tài liệu thành vector để tìm các đoạn gần nhau về mặt ngữ nghĩa \cite{reimers2019sbert,karpukhin2020dpr}. Cách tiếp cận này phù hợp với câu hỏi được diễn đạt khác với văn bản gốc. Tuy nhiên, nếu chỉ dùng tìm kiếm vector đơn lẻ, hệ thống vẫn có thể bỏ sót các thuật ngữ hiếm, tên lệnh hoặc ký hiệu đặc thù của Tin học. Ngoài ra, semantic search không tự giải quyết vấn đề sinh câu trả lời cuối cùng.

Nhóm thứ tư là các chatbot dùng LLM trực tiếp, chẳng hạn các nền tảng AI đa mục đích như ChatGPT, Gemini hoặc Claude. Các mô hình ngôn ngữ hiện đại được xây dựng trên nền tảng Transformer và đã được khảo sát rộng rãi trong bối cảnh giáo dục \cite{vaswani2017attention,wang2024llmeducation}. Ưu điểm của nhóm này là khả năng hội thoại tự nhiên, diễn giải linh hoạt và xử lý tốt nhiều yêu cầu phức tạp. Tuy nhiên, khi áp dụng trực tiếp vào SGK Tin học Việt Nam, các mô hình này thiếu ràng buộc về nguồn tài liệu, dễ trả lời vượt phạm vi chương trình hoặc sinh thông tin không có căn cứ trong SGK.

Nhóm thứ năm là các hệ thống chatbot RAG. RAG kết hợp truy xuất với sinh văn bản, cho phép hệ thống lấy tài liệu liên quan trước khi gọi LLM \cite{lewis2020rag}. Đây là hướng phù hợp hơn với dữ liệu SGK vì tri thức có thể được cập nhật bằng cách thay đổi kho tài liệu thay vì huấn luyện lại mô hình. Dù vậy, hiệu quả của RAG phụ thuộc mạnh vào chất lượng phân đoạn, vector nhúng, truy xuất, xếp hạng lại và cách xây dựng ngữ cảnh.

Nhóm thứ sáu là RAG lai (hybrid) hoặc RAG thích ứng (adaptive). Nhóm này mở rộng RAG cơ bản bằng cách kết hợp nhiều chiến lược tìm kiếm, lọc siêu dữ liệu theo phạm vi, viết lại truy vấn, truy xuất nhiều truy vấn và xếp hạng lại. Các thành phần như BM25, hợp nhất thứ hạng và xếp hạng lại bằng mô hình học sâu là các kỹ thuật thường gặp trong nhóm này \cite{robertson2009bm25,cormack2009rrf,nogueira2019passage}. Cách tiếp cận này phù hợp với dữ liệu SGK có cấu trúc phân cấp, nhưng đi kèm độ phức tạp cao hơn về thiết kế, ghi nhật ký và đánh giá.

\subsection{Bảng so sánh các nhóm giải pháp}

Để làm rõ mức độ phù hợp của từng nhóm giải pháp đối với bài toán chatbot giáo dục dựa trên SGK, Bảng~\ref{tab:solution-comparison-ch2} đối chiếu các nhóm giải pháp theo năm tiêu chí chính.

\begin{table}[H]
\centering
\small
\begin{tabular}{|p{0.19\textwidth}|p{0.16\textwidth}|p{0.16\textwidth}|p{0.17\textwidth}|p{0.16\textwidth}|}
\hline
\textbf{Nhóm giải pháp} & \textbf{Hiểu câu hỏi tự nhiên} & \textbf{Dùng dữ liệu riêng} & \textbf{Giảm hallucination} & \textbf{Phù hợp với SGK} \\ \hline
Dựa trên luật & Thấp & Trung bình & Cao nếu luật đủ & Thấp \\ \hline
FAQ/Tìm kiếm từ khóa & Thấp--Trung bình & Có & Trung bình & Trung bình \\ \hline
Tìm kiếm ngữ nghĩa & Trung bình & Có & Phụ thuộc tài liệu trả về & Trung bình--Cao \\ \hline
LLM trực tiếp & Cao & Thấp nếu không nạp context & Thấp & Thấp--Trung bình \\ \hline
RAG & Cao & Cao & Trung bình--Cao & Cao \\ \hline
RAG lai/thích ứng & Cao & Cao & Cao hơn nếu truy xuất tốt & Cao \\ \hline
\end{tabular}
\caption{So sánh các nhóm giải pháp cho bài toán chatbot giáo dục dựa trên SGK}
\label{tab:solution-comparison-ch2}
\end{table}

Bảng~\ref{tab:solution-comparison-ch2} cho thấy các nhóm giải pháp đơn giản như dựa trên luật hoặc FAQ có ưu điểm về khả năng kiểm soát trong phạm vi hẹp, nhưng khó mở rộng khi số lượng bài học và cách đặt câu hỏi tăng lên. LLM trực tiếp có khả năng xử lý ngôn ngữ tốt hơn, nhưng thiếu cơ chế gắn câu trả lời với dữ liệu SGK. RAG và RAG lai/thích ứng phù hợp hơn với bài toán của đồ án vì vừa tận dụng được năng lực sinh ngôn ngữ của LLM, vừa có thể kiểm soát nguồn tri thức thông qua retrieval.

\subsection{Nhận xét và định vị giải pháp đồ án}

Từ khảo sát trên, có thể rút ra ba nhận xét chính. Thứ nhất, các hệ thống dựa trên luật hoặc FAQ chưa đáp ứng tốt nhu cầu hỏi đáp tự nhiên của học sinh. Người học thường không đặt câu hỏi theo đúng mẫu cố định, do đó hệ thống cần khả năng hiểu ngữ nghĩa và xử lý nhiều cách diễn đạt khác nhau. Thứ hai, LLM trực tiếp tuy mạnh về ngôn ngữ nhưng không bảo đảm câu trả lời bám SGK. Với giáo dục phổ thông, một câu trả lời nghe hợp lý nhưng vượt phạm vi hoặc sai nguồn có thể gây tác động tiêu cực đến quá trình học. Thứ ba, RAG là hướng tiếp cận phù hợp, nhưng cần được thiết kế và đánh giá cẩn thận vì lỗi truy xuất có thể lan truyền sang bước sinh văn bản.

Khoảng trống mà đồ án hướng tới nằm ở giao điểm giữa ba yêu cầu: dữ liệu SGK tiếng Việt có cấu trúc phân cấp, câu hỏi tự nhiên từ người dùng, và nhu cầu đánh giá riêng từng node trong pipeline. Nhiều hệ thống hiện có chỉ đánh giá câu trả lời cuối cùng hoặc tập trung vào trải nghiệm sản phẩm, trong khi chưa phân tích rõ chất lượng của truy xuất, xếp hạng lại và sinh văn bản. Đồ án vì vậy định vị giải pháp theo hướng RAG lai/thích ứng, trong đó hệ thống có thể truy xuất theo phạm vi sách, lớp, chủ đề, xếp hạng lại ứng viên và xây dựng benchmark riêng để đo chất lượng truy xuất.

Đối với nhóm giải pháp thương mại đa mục đích, điểm khác biệt của đồ án không nằm ở việc xây dựng một LLM mới, mà nằm ở cách tổ chức tri thức SGK và kiểm soát ngữ cảnh đưa vào LLM. Đối với nhóm chatbot dịch vụ trong nước, điểm khác biệt nằm ở mục tiêu học thuật: hệ thống cần giải thích, truy xuất bài học, sinh học liệu và giữ ngữ cảnh sư phạm. Đối với các hệ thống nghiên cứu hoặc đồ án đi trước, điểm khác biệt nằm ở việc tách bộ truy xuất và bộ sinh văn bản thành các thành phần có thể quan sát, kiểm thử và đánh giá riêng.

\section{Cơ sở lý thuyết và hướng tiếp cận}

\subsection{Mô hình ngôn ngữ lớn và hạn chế khi dùng trực tiếp}

LLM là mô hình học máy có khả năng sinh văn bản tự nhiên dựa trên ngữ cảnh đầu vào. Kiến trúc Transformer là nền tảng quan trọng cho nhiều LLM hiện đại \cite{vaswani2017attention}. Trong bài toán giáo dục, LLM có thể diễn giải khái niệm, tóm tắt nội dung, chuyển đổi yêu cầu thành câu trả lời có cấu trúc và hỗ trợ sinh học liệu \cite{wang2024llmeducation}. Tuy nhiên, năng lực ngôn ngữ của LLM không đồng nghĩa với việc mô hình luôn có tri thức đúng theo SGK đang xét. Khi không được cung cấp ngữ cảnh từ tài liệu, mô hình có thể trả lời dựa trên tri thức tham số đã học, trong đó có thể chứa thông tin ngoài chương trình, thông tin không cập nhật hoặc thông tin không phù hợp với cách diễn đạt của SGK Việt Nam.

Hạn chế này đặc biệt quan trọng trong môi trường giáo dục. Một câu trả lời sai hoặc vượt cấp không chỉ là lỗi kỹ thuật, mà còn có thể làm người học hiểu sai kiến thức. Ngoài ra, cửa sổ ngữ cảnh (context window) của LLM có giới hạn, nên không thể đưa toàn bộ SGK vào câu lệnh cho mọi lượt hỏi. Vì vậy, đồ án không dùng LLM như nguồn tri thức duy nhất, mà đặt LLM sau lớp truy xuất để mô hình sinh câu trả lời dựa trên các chunk đã được truy xuất.

\subsection{Embedding và tìm kiếm ngữ nghĩa}

Vector nhúng (Embedding) là kỹ thuật biểu diễn văn bản thành vector số trong không gian nhiều chiều. Khi hai đoạn văn có nội dung gần nhau, vector của chúng thường có độ tương đồng cao. Trong hệ thống hỏi đáp dựa trên SGK, embedding cho phép đưa câu hỏi của người dùng và các chunk tài liệu vào cùng không gian biểu diễn, từ đó thực hiện tìm kiếm vector để tìm các đoạn có ý nghĩa gần với câu hỏi \cite{reimers2019sbert,karpukhin2020dpr}.

Tìm kiếm ngữ nghĩa có ưu điểm ở các truy vấn diễn đạt lại, dùng từ đồng nghĩa hoặc không trùng bề mặt với tài liệu. Tuy nhiên, tìm kiếm vector cũng có hạn chế khi gặp thuật ngữ hiếm, tên hàm, ký hiệu hoặc cụm từ kỹ thuật cần khớp chính xác. Vì lý do đó, hướng tiếp cận của đồ án không chỉ dựa vào vector nhúng đơn lẻ, mà kết hợp thêm tìm kiếm từ khóa, lọc theo siêu dữ liệu và xếp hạng lại. Các công thức cụ thể cho BM25, độ tương đồng cosine và RRF được trình bày trong Chương 3 khi mô tả cơ chế truy hồi ngữ cảnh.

\subsection{Retrieval-Augmented Generation}

Sinh văn bản tăng cường truy xuất (RAG) là kiến trúc kết hợp giữa truy xuất và sinh văn bản \cite{lewis2020rag}. Trong pha chuẩn bị dữ liệu, tài liệu nguồn được làm sạch, chia thành các phân đoạn, tạo vector nhúng và lưu vào chỉ mục kèm siêu dữ liệu. Trong pha truy vấn, hệ thống nhận câu hỏi, truy xuất các phân đoạn liên quan, có thể xếp hạng lại ứng viên, sau đó đưa ngữ cảnh vào câu lệnh để LLM sinh câu trả lời.

RAG phù hợp với bài toán SGK vì tri thức của hệ thống được quản lý thông qua kho tài liệu thay vì phụ thuộc hoàn toàn vào mô hình. Khi bổ sung hoặc cập nhật tài liệu, hệ thống có thể tái lập chỉ mục mà không cần huấn luyện lại LLM. Ngoài ra, các chunk được truy xuất có thể được ghi nhật ký, hiển thị hoặc dùng cho phân tích lỗi, giúp quá trình đánh giá minh bạch hơn so với việc chỉ quan sát câu trả lời cuối cùng.

Tuy nhiên, RAG không tự bảo đảm chất lượng đầu ra. Nếu tài liệu được phân đoạn không hợp lý, siêu dữ liệu sai, bộ truy xuất lấy nhầm bài hoặc ngữ cảnh đưa vào câu lệnh chứa nhiều nhiễu, LLM vẫn có thể sinh câu trả lời không đạt yêu cầu. Vì vậy, phương pháp của đồ án xem truy xuất là thành phần cần đánh giá định lượng riêng, còn việc sinh câu trả lời được đánh giá theo mức độ đúng, đủ, bám sát ngữ cảnh và rõ ràng.

\subsection{Adaptive retrieval, reranking và viết lại truy vấn}

Một truy vấn người dùng có thể thiếu thông tin về phạm vi, phụ thuộc vào hội thoại trước đó hoặc dùng cách diễn đạt không trùng với tài liệu. Kỹ thuật viết lại truy vấn giúp chuẩn hóa hoặc làm rõ truy vấn trước khi truy xuất. Trong khi đó, truy xuất thích ứng chọn chiến lược truy xuất tùy theo loại câu hỏi, phạm vi và tín hiệu ngữ cảnh, thay vì dùng một cấu hình cố định cho mọi trường hợp.

Xếp hạng lại (Reranking) là bước đánh giá lại danh sách candidate sau truy xuất ban đầu. Bộ truy xuất thường ưu tiên tốc độ và độ phủ, nên top-k ban đầu có thể chứa các phân đoạn chưa thật sự phù hợp. Bộ xếp hạng lại (Reranker), thường được triển khai bằng các mô hình học sâu như cross-encoder, nhận cặp truy vấn--tài liệu (query--document) và ước lượng lại mức độ liên quan \cite{nogueira2019passage}. Cách làm này có thể cải thiện thứ tự ngữ cảnh đưa vào LLM, nhưng cũng làm tăng độ trễ (latency) và chi phí tính toán. Vì vậy, trong Chương 4, kết quả thực nghiệm cần so sánh trước và sau khi xếp hạng lại để xác định lợi ích thực tế của bước này.

Trong đồ án, truy xuất thích ứng, viết lại truy vấn và xếp hạng lại được xem là các cơ chế bổ trợ cho lõi RAG. Chúng không thay thế mục tiêu chính là truy xuất đúng ngữ cảnh, mà giúp hệ thống xử lý tốt hơn các câu hỏi mơ hồ, câu hỏi có phạm vi cụ thể và câu hỏi được diễn đạt khác với văn bản SGK.

\subsection{Đa ý định và điều phối yêu cầu học tập}

Người dùng của hệ thống không chỉ đặt câu hỏi kiến thức, mà còn có thể yêu cầu tạo quiz, sinh slide, soạn giáo án hoặc kết hợp nhiều yêu cầu trong cùng một lượt hỏi. Vì vậy, hệ thống cần một lớp điều phối để phân tích ý định, xác định loại tác vụ (task type), lập kế hoạch hành động và gọi đúng dịch vụ. Lớp điều phối này giúp biến yêu cầu tự nhiên thành các hành động có cấu trúc, nhưng không phải trọng tâm đánh giá định lượng của đồ án.

Vai trò của việc điều hướng đa ý định (multi-intent routing) trong nghiên cứu này là bảo đảm yêu cầu được đưa vào đúng luồng xử lý. Với câu hỏi kiến thức, hệ thống cần ưu tiên truy xuất và sinh câu trả lời. Với yêu cầu sinh học liệu, hệ thống cần sử dụng ngữ cảnh từ SGK rồi chuyển sang luồng làm việc đa tác tử để xây dựng dàn ý, nội dung, phương tiện trực quan hoặc câu hỏi đánh giá. Cách tổ chức này cho phép các tác vụ khác nhau dùng chung nền tảng tri thức, đồng thời giữ cho truy xuất vẫn là thành phần trung tâm.

\subsection{Đánh giá hệ thống RAG}

Đánh giá hệ thống RAG cần tách hai câu hỏi: hệ thống có truy xuất đúng ngữ cảnh không, và LLM có sinh câu trả lời tốt từ ngữ cảnh đó không. Nếu chỉ đánh giá câu trả lời cuối cùng, rất khó xác định lỗi đến từ bộ truy xuất hay bộ sinh. Một câu trả lời sai có thể do ngữ cảnh sai, ngữ cảnh thiếu, câu lệnh chưa tốt hoặc LLM suy diễn quá mức. Ngược lại, một câu trả lời đúng không nhất thiết chứng minh bộ truy xuất hoạt động tốt, vì LLM có thể dựa vào tri thức sẵn có để trả lời.

Đối với truy xuất, đồ án sử dụng benchmark có nguồn kỳ vọng rõ ràng để tính các chỉ số như Hit@k, Recall@k, Precision@k, MRR và nDCG@k. Đây là các độ đo thường dùng trong đánh giá hệ thống truy hồi thông tin, trong đó nDCG được dùng để phản ánh chất lượng xếp hạng có xét vị trí kết quả \cite{thakur2021beir,jarvelin2002cumulated}. Đối với khâu sinh văn bản, hướng đánh giá phù hợp hơn là dùng rubric về độ đúng nội dung, độ đầy đủ, mức độ bám sát ngữ cảnh, độ rõ ràng và khả năng tránh hiện tượng ảo giác. Cách tách đánh giá này giúp quá trình phân tích lỗi có cơ sở hơn và phù hợp với mục tiêu nghiên cứu của đồ án.

\subsection{Khoảng trống và hướng tiếp cận của đồ án}

Từ các phân tích trên, có thể thấy khoảng trống chính của bài toán nằm ở việc xây dựng một hệ thống RAG cho dữ liệu SGK tiếng Việt có khả năng kiểm soát phạm vi và đánh giá được từng thành phần. Các giải pháp hiện có hoặc quá cứng nhắc khi dựa vào luật và từ khóa, hoặc quá mở khi dùng LLM trực tiếp, hoặc chưa đánh giá rõ chất lượng truy xuất trước khi kết luận về chất lượng câu trả lời.

Hướng tiếp cận của đồ án là thiết kế một luồng RAG cho SGK Tin học THPT, trong đó tài liệu được phân đoạn theo cấu trúc bài học, mỗi phân đoạn được gắn siêu dữ liệu phục vụ lọc phạm vi, truy xuất kết hợp tìm kiếm từ khóa và tìm kiếm vector, ứng viên được xếp hạng lại trước khi đưa vào khâu sinh câu trả lời, và các bước xử lý được ghi nhật ký để phục vụ đánh giá. Trên nền tảng đó, hệ thống mở rộng sang các tác vụ sinh học liệu bằng lớp điều phối và luồng làm việc đa tác tử, nhưng vẫn giữ nguyên nguyên tắc bám sát (grounding) trên nội dung SGK.

Phần tiếp theo của báo cáo sẽ trình bày chi tiết phương pháp đề xuất. Cụ thể, Chương 3 mô tả kiến trúc tổng thể, nền tảng tri thức, lõi RAG, tầng điều phối tác tử, tầng sinh nội dung đa tác tử và phương pháp đánh giá. Những nội dung nền tảng trong chương này là cơ sở để giải thích các quyết định thiết kế và thực nghiệm ở các chương sau.

\section{Kết chương}

Chương này đã xác định phạm vi nghiên cứu của đồ án là hệ thống hỏi đáp dựa trên SGK Tin học THPT theo hướng document-grounded QA. Chương cũng đã phân tích đặc thù của dữ liệu SGK, làm rõ các thách thức về phạm vi, ngôn ngữ tự nhiên, thuật ngữ Tin học và vai trò của bộ truy xuất trong chất lượng câu trả lời.

Thông qua khảo sát các nhóm giải pháp liên quan, chương đã chỉ ra rằng RAG và RAG lai/thích ứng là hướng tiếp cận phù hợp hơn so với hệ thống dựa trên luật, FAQ hoặc LLM trực tiếp khi yêu cầu câu trả lời phải bám sát SGK. Các kiến thức nền tảng về LLM, vector nhúng, RAG, viết lại truy vấn, xếp hạng lại, điều hướng đa ý định và đánh giá hệ thống RAG tạo cơ sở cho Chương 3, nơi phương pháp đề xuất được trình bày chi tiết hơn về kiến trúc, các nút (node) xử lý và phương pháp đánh giá.

\end{document}
